\chapter*{} % Unnumbered chapter
\addcontentsline{toc}{chapter}{Part V: Foundation Model Prior in SLAM Pipeline } % Add to ToC
\begin{center}
    \vspace*{3.5cm} 
{\Huge \textbf{Part V} \\ [1cm]}
{\LARGE \textbf{Foundation Model Prior in SLAM Pipeline }}
\end{center}


\begin{figure}[H]  % Force the figure to stay HERE
    \centering
    \includegraphics[width=0.95\linewidth]{figs/ch0/wild4d-slam.png}
    \caption{Chapter Overview}
    \label{fig:en}
\end{figure}



In the previous chapter \ref{chapter:superodom},\ref{chapter:superodom2.0}, we demonstrated that multi-sensor fusion can enhance SLAM robustness. Yet, adding additional sensors is often impractical—it raises hardware costs, system complexity, and computational load. A key open question is how to deliver reliable performance without “simply adding more sensors.” We address this challenge by embedding \textbf{foundation-model priors} within a SLAM pipeline. We focus on visual–inertial sensor setup, the most widely deployed and lightweight configuration on contemporary robotic platforms. In this chapter, we will focus on how to achieve real-time state estimation and 4D mapping in any environment (dynamic, dark and etc) with minimal sensor setups.

